% !PRINTED 624
\weberpage{624}
% Running head "624  Vierundzwanzigster Abschnitt.  § 171." is page furniture and
% is not transcribed as body text; only the printed number is recorded above.
%<<CONT
% (This page's own footnote apparatus begins mid-sentence, continuing the
% footnote 1) that was left open at the bottom of the previous page: "...Clebsch,
% Über die" + "Anwendung der Abelschen Funktionen..." joins into one sentence.
% The main-text continuation below is a fresh section, not a continued sentence.)

\begin{center}
\textbf{§ 171. Definition of algebraic functions.}
\end{center}

A variable $\theta$ is called an algebraic function of an independent
variable $z$ when the two variables are connected with one another by an
algebraic equation
\weq{1}{F(\theta,z)=0}
In this, $F$ denotes an expression of the form
\weq{2}{F(\theta,z)=a_0\theta^n+a_1\theta^{n-1}+\cdots+a_{n-1}\theta+a_n,}
whose coefficients $a_0, a_1, \ldots, a_n$ are integral rational functions
of $z$ without a common divisor. About the constant coefficients in these
expressions we make no further assumption than that they are real or
complex numbers\textsuperscript{1)}.

We assume the function $F(\theta,z)$
to be irreducible in the sense that it shall not split into factors which
are themselves rational functions of $\theta$ and $z$.

Every integral
function $G(\theta,z)$ can, according to Bd.\,I, § 20, be decomposed in
only one way into irreducible factors (apart from constant factors). We
say that the function $G(\theta,z)$
%CONT>>
% body text is cut off here by the page break; the sentence continues on the
% next page, which is not included in this scan set.

\smallskip
\noindent\rule{5em}{0.4pt}

\noindent Anwendung der \gesperrt{Abelschen} Funktionen in der Geometrie, Crelles
Journ., Bd.\,63, 1864. \gesperrt{Clebsch} and \gesperrt{Gordan},
Theorie der \gesperrt{Abelschen} Funktionen, Leipzig 1866. \gesperrt{Brill} and
\gesperrt{Noether}, Über die algebraischen Funktionen und ihre
Anwendung in der Geometrie, Mathematische Annalen VII, 1874.
\gesperrt{Brill} and \gesperrt{Noether}, Die Entwickelung der Theorie
der algebraischen Funktionen, Bericht der Deutschen
Mathematiker-Vereinigung 1894. \gesperrt{Appell} et \gesperrt{Goursat},
Théorie des fonctions algébriques, Paris 1895. On the
number-theoretic methods: \gesperrt{Kronecker}, Über die
Diskriminante algebraischer Funktionen, Crelles Journ., Bd.\,91, 1881,
and: Festschrift zu \gesperrt{Kummers} Doktor-Jubiläum, Crelles
Journ., Bd.\,92, 1881. \gesperrt{Dedekind} and \gesperrt{Weber},
Theorie der algebraischen Funktionen einer Veränderlichen, Crelles
Journ., Bd.\,92, 1879. \gesperrt{Hensel} and \gesperrt{Landsberg},
Theorie der algebraischen Funktionen einer Variablen, Leipzig 1902.
The form of the theory presented in what follows rests upon the concept
of functionals introduced in Bd.\,II, § 153, which in this form had
already been applied, in a Strassburg dissertation prompted by
\gesperrt{Wellstein} and written by \gesperrt{Rehfeld} (1906), to the
treatment of a special problem of the theory of functions.

\smallskip
\noindent\textsuperscript{1)}~No gap would remain anywhere if we wished
to restrict ourselves here to \gesperrt{algebraic numbers}. If one regards
$\theta, z$ as Cartesian coordinates in the plane, then equation (1)
represents an algebraic curve, and one can apply the theory of these
curves, as \gesperrt{Clebsch} and \gesperrt{Gordan} and \gesperrt{Brill}
and \gesperrt{Noether} have done. To be sure, only the real values of
these variables have a really intuitive meaning here.
